\section{A total eventual-sign worker with explicit witnesses}
\label{sec:tails}

\subsection{Why eventual goals deserve their own certificate}

A positivity certificate on $[a,\infty)$ answers a different question from
``there is some $a$ beyond which the function is positive.'' The second query
should not be handled by guessing a large number until a bounded search happens
to succeed. For \eqref{eq:fragment}, a direct exact algorithm constructs a
cutoff for every nonzero canonical input.

This worker is particularly useful for existential goals. It returns a rational
witness $X$ together with a proof obligation that is uniform over every
$x\ge X$. It can also supply the tail component of a proof whose bounded
component is discharged by another worker. A tail receipt alone must never be
silently promoted to a proof over the original entire domain.

\subsection{An explicit domination theorem}

Let $f\ne0$, let $\lambda_*$ be its largest rate, and let $d$ be the degree of
$P_{\lambda_*}$. Multiply the target by a sign $\sigma\in\{-1,1\}$ so that the
leading coefficient of $\sigma P_{\lambda_*}$ is $a>0$. Define
\begin{align*}
 L&=\sum_{k<d}|[x^k]P_{\lambda_*}|,\\
 C_\lambda&=\sum_k|[x^k]P_\lambda|,\qquad
 d_\lambda=\deg P_\lambda,\qquad
 \delta_\lambda=\lambda_*-\lambda>0,\\
 k_\lambda&=\max(d_\lambda-d,0)+1,\\
 B&=\sum_{\lambda<\lambda_*}
       \frac{C_\lambda k_\lambda!}{\delta_\lambda^{k_\lambda}}.
\end{align*}
Empty sums are zero. All these quantities are rational except for the integer
degrees and orders.

\begin{theorem}[Exact eventual-sign certificate]
\label{thm:tail}
For
\begin{equation}
 X=\max\left(1,\frac{2L}{a},\frac{4B}{a}\right),
 \label{eq:cutoff}
\end{equation}
one has
\begin{equation}
 \sigma f(x)\ \ge\ \frac{a}{4}x^d e^{\lambda_*x}>0
 \qquad\text{for every }x\ge X.
 \label{eq:tail-lower}
\end{equation}
\end{theorem}
\begin{proof}
For $x\ge1$, the lower-degree terms of the leading block have total magnitude
at most $Lx^{d-1}$; when $d=0$, there are no such terms and $L=0$. Hence
\[
 \sigma P_{\lambda_*}(x)
 \ge ax^d-Lx^{d-1}\ge\frac{a}{2}x^d
\]
when $x\ge2L/a$.

For a lower-rate block, $|P_\lambda(x)|\le C_\lambda x^{d_\lambda}$.
The exponential series, or the Taylor-remainder ladder of Section
\ref{sec:ladders}, gives $e^t\ge t^k/k!$ for $t\ge0$ and $k\in\N$.
Since $\delta_\lambda x>0$,
\begin{align*}
 |P_\lambda(x)|e^{-\delta_\lambda x}
 &\le\frac{C_\lambda k_\lambda!}
             {\delta_\lambda^{k_\lambda}}
       x^{d_\lambda-k_\lambda}\\
 &\le\frac{C_\lambda k_\lambda!}
             {\delta_\lambda^{k_\lambda}}
       x^{d-1}.
\end{align*}
The last step uses $d_\lambda-k_\lambda\le d-1$ and $x\ge1$.
Summing yields a lower-rate contribution of magnitude at most
$Bx^{d-1}\le(a/4)x^d$ when $x\ge4B/a$. Factoring out
$e^{\lambda_*x}>0$ now gives \eqref{eq:tail-lower}. Strictness follows from
$a>0$ and $x\ge1$.
\end{proof}

\begin{corollary}[Eventual sign and identity are decidable in the exact fragment]
The eventual sign of a nonzero canonical exponential polynomial is the sign
of the leading coefficient in its largest-rate block. The zero canonical
object is identically zero. Conversely, a nonzero canonical object cannot
represent the zero function, by Theorem~\ref{thm:tail}.
\end{corollary}

This is a genuine completeness statement for the eventual-sign problem in the
specified rational fragment, subject only to the practical limits of exact
computation. It does not decide where the sign changes, whether the function
is nonnegative on an earlier compact interval, or whether its minimum is zero.
Those are different problems.

\subsection{The search and the checker}

The producer extracts the largest rate, leading degree and sign, computes the
norms above, selects the displayed $k_\lambda$, and emits $X$ with those orders.
The checker independently requires
\[
 X\ge1,\quad a>0,\quad aX\ge2L,\quad aX\ge4B,
\]
and verifies for every lower block that its supplied integer $k_\lambda$ is
positive and satisfies $k_\lambda\ge d_\lambda-d+1$.
It does not require the producer's exact order choice or its exact minimum
cutoff. A different search can choose larger orders or a larger cutoff without
changing the acceptance theorem.

The deliberately elementary cutoff may be quite conservative. For
\[
 f(x)=e^x-1000x^5
\]
one has $a=1$, $L=0$, $k=6$, $B=1000\cdot6!=720000$; thus $X=2880000$ is
an immediate certificate. This is not a useful numerical estimate of the last
zero. It is a total, exact fallback that makes no root-isolation assumption.

A cheap refinement tries the native ladder at increasing anchors. On this
example the tested sequence $0,1,2,4,8,16,32$ first obtains a certificate at
$32$. That certificate proves strict positivity for all $x\ge32$ using exact
rational enclosures of its initial values. The experiment establishes neither
that $32$ is the optimal cutoff nor that this anchor search always stays cheap.
Its role is to improve a witness whose existence is already established by a
different total algorithm.

\subsection{Why sufficiently large native anchors work mathematically}

There is a useful bridge between the two workers. Suppose the eventual leading
sign is positive and use the native ascending factors. Every nonzero prefix
function also has positive eventual leading sign. Before factors for the
largest rate are reached, each lower factor $(\D-\rho)$ multiplies the leading
coefficient of the largest-rate block by the positive quantity
$\lambda_*-\rho$. Once the largest-rate factors are reached, all lower-rate
blocks are already dead; the remaining factors differentiate its polynomial,
preserving a positive leading coefficient until the function vanishes.

There are finitely many prefixes. Apply Theorem~\ref{thm:tail} to each nonzero
one and choose an anchor larger than all resulting cutoffs. Every initial
coordinate is then positive, and the native ladder succeeds. Thus a sufficiently
large anchor exists without any auxiliary factors.

This argument distinguishes mathematical termination from the delivered bounded
search. A fair search over unbounded anchors and increasing enclosure precision
can eventually establish all the strictly positive initial values. The actual
experiment tries a finite sequence at fixed precision, with the explicit tail
certificate as its fallback. It makes no universal runtime claim.

\subsection{Quantitative asymptotics available from the same proof}

The estimates also give, for $x\ge1$,
\begin{equation}
 \left|\sigma e^{-\lambda_*x}x^{-d}f(x)-a\right|
 \le\frac{L+B}{x}.
 \label{eq:relative-modulus}
\end{equation}
Consequently, for any rational $\eta>0$, the rational witness
\[
 X_\eta=\max\left(1,\frac{L+B}{a\eta}\right)
\]
ensures relative error at most $\eta$ with respect to the leading term.
The proof is the same triangle-inequality estimate, without splitting the
budget into halves and quarters.

This corollary suggests an additional Forge contract:
\par\noindent\begin{minipage}{\linewidth}
\begin{lstlisting}
leadingTermWithModulus(f, eta) ->
    (sign, rate, degree, coefficient, cutoff, proof)
\end{lstlisting}
\end{minipage}\par
Such a result can feed ratio comparisons, limits, and eventual denominator
nonvanishing. The archive implements the eventual-sign contract, not this
separate modulus API; the mathematical corollary is stated explicitly so that
the next implementation does not need to discover another method.

If the source variable is restricted to naturals, replacing a rational cutoff
by its ceiling gives an integer witness after a proved coercion step. If the
source is a logarithmic chart $x=e^t$, a rational cutoff in $t$ corresponds to
the real witness $e^X$ in $x$. The witness need not itself be rational. These
transport steps belong to source reconstruction, not to rational rounding in
the search process.

\subsection{Combining a tail with a bounded-domain proof}

For a requested ray $[a,\infty)$ and a tail cutoff $X\ge a$, produce two
separate subgoals:
\[
 \forall x\in[a,X],\ f(x)\ge0,
 \qquad
 \forall x\ge X,\ f(x)\ge0.
\]
The source proof uses the elementary dichotomy $x\le X$ or $X\le x$.
The shared endpoint must be covered; there is no half-open gap. If the target
is strict, both component contracts must carry enough strict information.

This decomposition is especially attractive for a function that decreases
initially and then becomes positive. It is not by itself a universal
nonnegativity procedure: a compact component containing a zero can still defeat
a simple interval worker. The derivative and factorisation routes are needed
precisely at those contacts.
